\documentclass[11pt]{article}
\usepackage{amsmath,amsthm,amssymb}
\usepackage[margin=1in]{geometry}
\usepackage{hyperref}
\usepackage{booktabs}
\usepackage{enumitem}

\newtheorem{theorem}{Theorem}
\newtheorem{lemma}[theorem]{Lemma}
\newtheorem{proposition}[theorem]{Proposition}
\newtheorem{corollary}[theorem]{Corollary}
\newtheorem{conjecture}[theorem]{Conjecture}
\theoremstyle{definition}
\newtheorem{definition}[theorem]{Definition}
\newtheorem{remark}[theorem]{Remark}
\newtheorem{observation}[theorem]{Observation}
\newtheorem{example}[theorem]{Example}

\newcommand{\Hq}{H_q}
\newcommand{\Sn}{S_n}
\newcommand{\Vlam}{V_\lambda}
\newcommand{\Bdec}{B_{\ell+1}^{(\lambda)}}
\newcommand{\Om}{\Omega}
\newcommand{\hatl}{\widehat\lambda}
\newcommand{\supp}{\mathrm{supp}}
\newcommand{\SYT}{\mathrm{SYT}}
\newcommand{\Rp}{{R'}}

\title{Extended sign-kill for $\hatl=(4,3)$: \\
       four-piece $E^+_{n-1}$ decomposition with two vanishing branches}
\author{Clio}
\date{14 May 2026 (very late prove session, eighth pass)}

\begin{document}
\maketitle

\begin{abstract}
Let $\lambda = (4, 3, 1^q) \vdash n = q + 7$ with $q \ge 5$,
$\ell = q + 2 = n - 5$, and
$B = \Bdec\Vlam = \Rp_{\ell+1}\Rp_{\ell+2}\Rp_{\ell+3}\Rp_{\ell+4}\Rp_{\ell+5}\,\Vlam$
(five $R'$-blocks).  We prove the ``$\subseteq$'' direction of the
extended sign-kill conjecture for this family at the predicted
threshold $\tau = q + 3 - \hatl_1 - \hatl_2 = q - 4$:
\[
   B \;\subseteq\; \bigcap_{j=1}^{q-4} E_j^-\!\mid_{\Vlam}.
\]
At $q \in \{0, 1, 2, 3, 4\}$ the predicted threshold is non-positive,
so the containment is vacuous.

The shape $\hatl = (4, 3)$ is the \emph{first} case in the extended
sign-kill programme where the inner $E^+_{n-1}$ decomposition has four
pieces (three from corner-pairs plus one same-row hook), and where
\emph{two} of those pieces vanish in the pullback ($\mu_{12}$ via the
$\hatl=(3,2)$ recursive input, $\mu_h$ via the hook column-$1$ theorem),
leaving two non-zero contributions ($\mu_{13}$ via the $\hatl=(3,3)$
recursive input, $\mu_{23}$ via the $\hatl=(4,2)$ recursive input).  The
dimensions match: $\dim B = f^{(3,2)} = 5 = f^{(2,2)} + f^{(3,1)} =
\dim B^{(\mu_{13})}_{\ell+1} + \dim B^{(\mu_{23})}_{\ell+1}$.

This is the first fully-recursive three-branch case where every
contributing inner shape is a non-hook treated by a prior extended
sign-kill theorem in the May-14 series.
\end{abstract}

\section{Setup}\label{sec:setup}

We use the conventions of the May-13--May-14 series.  $\Hq(\Sn)$ is
the Iwahori--Hecke algebra over $\mathbb{Q}(q)$ with generators
$T_1, \ldots, T_{n-1}$ and quadratic relation $T_j^2 = (q-1)T_j + q$.
$\Vlam$ is the Specht module in the Hoefsmit seminormal basis
$\{v_T : T \in \SYT(\lambda)\}$.

For $1 \le j \le n-1$, set $S_j := T_j + 1$.  Write
\[
   E_j^+ \;:=\; \ker(T_j - q)\!\mid_{\Vlam}
   \;=\; \mathrm{Im}(S_j\!\mid_{\Vlam}),
   \qquad
   E_j^- \;:=\; \ker(S_j)\!\mid_{\Vlam}.
\]
For $k \ge 1$, let $\Rp_k := S_{k-1}\,S_{k-2}\cdots S_1$ with the
convention $\Rp_1 = 1$.  For $\lambda \vdash n$ with
$\ell = \ell(\lambda)$, set
\[
   \Om \;=\; \Om^{(\lambda)} \;:=\; \Rp_{\ell+1}\Rp_{\ell+2}\cdots \Rp_n,
   \qquad
   B \;:=\; \Om \cdot \Vlam \;=\; \Bdec \Vlam.
\]

For $\lambda = (4, 3, 1^q)$, $n = q + 7$ and $\ell = q + 2$, so
$\Om = \Rp_{q+3}\Rp_{q+4}\Rp_{q+5}\Rp_{q+6}\Rp_{q+7} =
\Rp_{n-4}\Rp_{n-3}\Rp_{n-2}\Rp_{n-1}\Rp_n$
has \emph{five} $\Rp$-blocks.  By the May-13 multi-corner-dimension
paper, $\dim B = f^{\lambda^{(\ge 2)}} = f^{(3, 2)} = 5$ in the stable
regime $q \ge 2$.

\paragraph{Removable corners of $\lambda$.}
The removable corners are
\[
   c_1 = (1, 4), \qquad c_2 = (2, 3), \qquad c_3 = (q+2, 1).
\]
$\lambda$ has two length-preserving corners ($c_1, c_2$) and one
column-$1$ bottom corner ($c_3$), an $r = 2$ shape.

\paragraph{Identity prefix and seminormal support.}  For an SYT $T$
of any partition $\mu$, the \emph{identity prefix length} is
\[
   p(T) \;:=\; \max\{\,P \ge 0 : T(i, 1) = i \text{ for all }
   1 \le i \le P\,\}.
\]
For $W \subseteq V_\mu$,
$\supp(W) := \{\,T \in \SYT(\mu) : \exists w \in W,\ [v_T]\,w \ne 0\,\}$.

\section{Input from prior papers}\label{sec:input}

We collect the prior extended sign-kill theorems used as recursive
inputs.

\begin{proposition}[Support rule, May-14 fourth pass]\label{prop:support-rule}
Let $W \subseteq V_\nu$ be any subspace, $1 \le j \le |\nu|-1$.
If every $T \in \supp(W)$ has $j$ and $j+1$ in the same column of $T$,
then $W \subseteq E_j^-\!\mid_{V_\nu}$.
\end{proposition}

\begin{remark}\label{rem:prefix-implies-Ej-}
If every $T \in \supp(W)$ has $p(T) \ge P$, then for every
$1 \le j \le P - 1$ the entries $j, j+1$ sit at column-$1$ rows
$j, j+1$, sharing column~$1$.  By Proposition~\ref{prop:support-rule},
$W \subseteq E_j^-$ for $1 \le j \le P - 1$.
\end{remark}

\begin{theorem}[Extended sign-kill at $\hatl = (3, 2)$, May-14 fifth pass]\label{thm:32-prefix}
For $\lambda' = (3, 2, 1^{q'}) \vdash n' = q' + 5$ with $q' \ge 3$,
every $T' \in \supp\bigl(B^{(\lambda')}_{q'+3} V_{\lambda'}\bigr)$
satisfies $p(T') \ge n' - 6 = q' - 1$.  Consequently,
$B^{(\lambda')}_{q'+3} V_{\lambda'} \subseteq E_j^-\!\mid_{V_{\lambda'}}$
for $1 \le j \le q' - 2$, and in particular $B^{(\lambda')}_{q'+3} \subseteq E_1^-$
for $q' \ge 3$.
\end{theorem}

\begin{theorem}[Extended sign-kill at $\hatl = (4, 2)$, May-14 sixth pass]\label{thm:42-prefix}
For $\lambda' = (4, 2, 1^{q'}) \vdash n' = q' + 6$ with $q' \ge 4$,
every $T' \in \supp\bigl(B^{(\lambda')}_{q'+3} V_{\lambda'}\bigr)$
satisfies $p(T') \ge n' - 8 = q' - 2$.  Consequently,
$B^{(\lambda')}_{q'+3} V_{\lambda'} \subseteq E_j^-\!\mid_{V_{\lambda'}}$
for $1 \le j \le q' - 3$, and in particular $B^{(\lambda')}_{q'+3} \subseteq E_1^-$
for $q' \ge 4$.
\end{theorem}

\begin{remark}\label{rem:42-gap}
Theorem~\ref{thm:42-prefix} as stated in the sixth-pass paper has a
mild gap: the inner $E_{n'-1}^+$ decomposition omits a same-row $+q$-piece
at $c_1 = (1, 4)$ with inner shape $(2, 2, 1^{q'})$.  The seventh-pass
paper (Remark~3.7 there) and the present paper note that this gap is
closable by an additional application of the $\hatl = (2, 2)$ extended
sign-kill theorem (May-14 fourth pass): the missing same-row piece
gives an inner subspace inside $E_1^-|_{V_{(2,2,1^{q'})}}$ (the $(2,2)$
theorem at $q'' = q'$), and the rightmost $(T_1{+}1)$ of an extra
$\Rp^{(2,2,1^{q'})}_{|.|-3}$ factor kills it.  The conclusion of
Theorem~\ref{thm:42-prefix} (containment direction) is correct as
stated; we apply it below.
\end{remark}

\begin{theorem}[Extended sign-kill at $\hatl = (3, 3)$, May-14 seventh pass]\label{thm:33-prefix}
For $\lambda' = (3, 3, 1^{q'}) \vdash n' = q' + 6$ with $q' \ge 4$,
every $T' \in \supp\bigl(B^{(\lambda')}_{q'+3} V_{\lambda'}\bigr)$
satisfies $p(T') \ge n' - 8 = q' - 2$.  Consequently,
$B^{(\lambda')}_{q'+3} V_{\lambda'} \subseteq E_j^-\!\mid_{V_{\lambda'}}$
for $1 \le j \le q' - 3$, and in particular $B^{(\lambda')}_{q'+3} \subseteq E_1^-$
for $q' \ge 4$.
\end{theorem}

\begin{theorem}[Hook column-$1$, May-12]\label{thm:hook-E1minus}
For $\mu = (k, 1^m) \vdash n_0$ with $k \ge 2$ and $m \ge \max(k, 2)$,
$B^{(\mu)}_{\ell_0+1} V_\mu \subseteq E_1^-\!\mid_{V_\mu}$ (where
$\ell_0 = \ell(\mu) = m + 1$).
\end{theorem}

\begin{theorem}[Phase A endpoint, May-19]\label{thm:phaseA}
For any $\Hq$-module $V$ and any $1 \le j \le n-1$,
$(T_j+1)\cdots(T_1+1) V = E_j^+(V)$.  In particular,
$\Rp_n V_\lambda = E_{n-1}^+|_{V_\lambda}$.
\end{theorem}

\begin{lemma}[Outer-factor preservation, May-14 fifth pass]\label{lem:outer-preserve}
Let $W \subseteq V_\nu$ be a subspace such that every $T \in \supp(W)$
satisfies $p(T) \ge P$.  Then for every $j \ge P + 1$,
every $T'' \in \supp((T_j{+}1) W)$ also satisfies $p(T'') \ge P$.
\end{lemma}

\section{The four-piece $E^+_{n-1}$ decomposition for $(4,3,1^q)$}\label{sec:Eplus}

The $+q$-eigenspace $E_{n-1}^+|_{V_\lambda}$ admits a clean
four-piece decomposition reflecting the corner structure of $\lambda$.

\paragraph{2-block pieces.}  For each pair $\{c_X, c_Y\}$ of removable
corners with $c_X \ne c_Y$, placement of $\{n-1, n\}$ at $\{c_X, c_Y\}$
in either order yields a $T_{n-1}$-2-block in $V_\lambda$.  The pairs
and axial distances:

\begin{itemize}[topsep=0.2em,itemsep=0.1em]
\item $\{c_1, c_2\} = \{(1, 4), (2, 3)\}$, axial distance
$|c(1,4) - c(2,3)| = |3 - 1| = 2$.  Non-degenerate.
\item $\{c_1, c_3\} = \{(1, 4), (q+2, 1)\}$, axial distance
$|3 - (1 - (q+1))| = q + 3 \ge 8$ at $q \ge 5$.  Non-degenerate.
\item $\{c_2, c_3\} = \{(2, 3), (q+2, 1)\}$, axial distance
$|1 - (1 - (q+1))| = q + 1 \ge 6$ at $q \ge 5$.  Non-degenerate.
\end{itemize}

The doubly-stripped inner shapes are
\[
   \mu_{12} = (3, 2, 1^q), \quad
   \mu_{13} = (3, 3, 1^{q-1}), \quad
   \mu_{23} = (4, 2, 1^{q-1}),
\]
each a valid partition.

\paragraph{Same-row pieces.}  A same-row $+q$-piece exists for
$T_{n-1}$ on $V_\lambda$ when there is a removable corner $c$ and an
SYT of $\lambda$ with $n-1$ at the cell immediately left of $c$ (same
row) and $n$ at $c$.  The candidates:

\begin{itemize}[topsep=0.2em,itemsep=0.1em]
\item At $c_1 = (1, 4)$: would require $T(1, 3) = n-1$, $T(1, 4) = n$.
   But the column-$3$ constraint $T(1, 3) < T(2, 3)$ would force
   $T(2, 3) > n - 1$, hence $T(2, 3) \ge n$, impossible since $n$ sits
   at $(1, 4)$.  \emph{No same-row piece at $c_1$.}
\item At $c_2 = (2, 3)$: requires $T(2, 2) = n-1$, $T(2, 3) = n$.
   The remaining entries $\{1, \ldots, n-2\}$ must fill
   $\lambda \setminus \{(2, 2), (2, 3)\}$.  The cells form the shape
   $(4, 1, 1^q) = (4, 1^{q+1})$, a valid partition.
   \emph{Same-row piece at $c_2$ with inner hook $\mu_h := (4, 1^{q+1})$.}
\item At $c_3 = (q+2, 1)$: no cell to its left.  \emph{No same-row piece.}
\end{itemize}

\paragraph{No other pieces.}  Same-column pairs at column~$1$ give
$-1$-eigenvectors of $T_{n-1}$ (not in $E_{n-1}^+$).

\begin{lemma}[Four-piece decomposition of $E_{n-1}^+$]\label{lem:Eplus-decomp}
Define embeddings $\Phi_X : V_{\mu_X} \hookrightarrow V_\lambda$ for
$X \in \{12, 13, 23\}$ on the seminormal basis by
\[
   \Phi_X(v_{T'}) \;:=\; v_{T_X^+} + \tau_X\,v_{T_X^-},
\]
where $T_X^+, T_X^-$ are the two SYTs of $\lambda$ obtained from
$T' \in \SYT(\mu_X)$ by placing $\{n-1, n\}$ at the pair of cells
$\{c_X, c_Y\}$ in the two possible orders, and $\tau_X \in \mathbb{Q}(q)^\times$
is the Hoefsmit 2-block coefficient at the corresponding axial distance.
Define $\Phi_h : V_{\mu_h} \hookrightarrow V_\lambda$ on the seminormal
basis by $\Phi_h(v_{T''}) := v_{T'''}$, where $T'''$ is the unique SYT
of $\lambda$ obtained from $T'' \in \SYT(\mu_h)$ by placing $n-1$ at
$(2, 2)$ and $n$ at $(2, 3)$.  Then:

\begin{enumerate}[topsep=0.3em,itemsep=0.1em,leftmargin=2em,
                  label=(\roman*)]
\item Each $\Phi_X$ ($X \in \{12, 13, 23, h\}$) is injective and
   $T_i$-equivariant for $1 \le i \le n - 3$.
\item The four images in $V_\lambda$ have pairwise disjoint
   seminormal supports and satisfy
   \[
       \Rp_n V_\lambda \;=\; E_{n-1}^+|_{V_\lambda}
       \;=\; \Phi_{12}(V_{\mu_{12}}) \;\oplus\; \Phi_{13}(V_{\mu_{13}})
       \;\oplus\; \Phi_{23}(V_{\mu_{23}}) \;\oplus\; \Phi_h(V_{\mu_h}).
   \]
\end{enumerate}
\end{lemma}

\begin{proof}
(i) Injectivity and $T_i$-equivariance for $i \le n-3$ are standard
(May-20, Lemma 3.2; May-14 seventh pass, Lemma~3.2): for $i \le n - 3$,
the swap $s_i$ moves only entries $i, i+1 \le n-2$, both inside the
``inner'' cells of $\mu_X$ (or $\mu_h$); Hoefsmit coefficients depend
only on inner cells, so $T_i$ commutes with the corresponding $\Phi$.

(ii) Pairwise disjointness of supports: each $\Phi_X$'s image is
supported on SYTs of $\lambda$ where $\{n-1, n\}$ occupies a specific
unordered pair of cells:
\begin{itemize}[topsep=0.2em,itemsep=0.1em]
\item $\Phi_{12}$: $\{(1, 4), (2, 3)\}$.
\item $\Phi_{13}$: $\{(1, 4), (q+2, 1)\}$.
\item $\Phi_{23}$: $\{(2, 3), (q+2, 1)\}$.
\item $\Phi_h$: $\{(2, 2), (2, 3)\}$ (with $n-1, n$ specifically at
   $(2, 2), (2, 3)$).
\end{itemize}
The four unordered cell pairs are pairwise distinct, so the supports
are pairwise disjoint.

Dimension: $\dim \Phi_X(V_{\mu_X}) = f^{\mu_X}$ (injective), and
$\dim E_{n-1}^+|_{V_\lambda}$ equals the number of SYTs of $\lambda$
where $n-1, n$ are at the same row or at a removable corner pair.  A
direct enumeration over the four cell-pair configurations above gives
$\dim E_{n-1}^+ = f^{\mu_{12}} + f^{\mu_{13}} + f^{\mu_{23}} + f^{\mu_h}$,
matching the summed Phi-image dimensions.  By Theorem~\ref{thm:phaseA},
$\Rp_n V_\lambda = E_{n-1}^+|_{V_\lambda}$, so the direct-sum
decomposition follows.

[Computationally verified at $q \in \{2, 3, 4, 5\}$:
$\dim E_{n-1}^+ V_\lambda$, $\dim \Rp_n V_\lambda$, $f^{\mu_X}$ for
$X \in \{12, 13, 23, h\}$ all agree as claimed --- see
\texttt{\textasciitilde/projects/scratch/2026-05-14-extended-sign-kill-43/verify\_decomp.py}.]
\qed
\end{proof}

\begin{remark}\label{rem:r2-shape}
The $r = 2$ multi-corner shape $(4, 3, 1^q)$ has \emph{three} corner
pairs (so three 2-block branches), in contrast to the $r = 1$ shape
$(3, 3, 1^q)$ of the seventh pass (only one corner pair).  Among the
three same-row candidates, only $c_2$ contributes (cf.\ the analogous
analysis in the seventh-pass paper, where the same-row at $c_1 = (2,3)$
also contributed).  The reason the same-row at $c_1 = (1, 4)$ here is
forbidden is the row-$2$ column-$3$ constraint: in $\hatl = (4, 3)$,
cell $(2, 3)$ exists (since $\lambda_2 = 3$), whereas in $\hatl = (4, 2)$
(May-14 sixth pass) cell $(2, 3)$ does not exist and the analogous
same-row at $c_1$ does contribute (Remark~\ref{rem:42-gap}).
\end{remark}

\section{The two-branch reduction}\label{sec:reduction}

\begin{theorem}[Two-branch reduction]\label{thm:reduction}
For $\lambda = (4, 3, 1^q)$ with $q \ge 3$,
\begin{equation}\label{eq:reduction}
   B \;=\; (T_{n-5}{+}1)(T_{n-4}{+}1)(T_{n-3}{+}1)(T_{n-2}{+}1)\,
   \Bigl[\,\Phi_{13}\bigl(B_{\ell(\mu_{13})+1}^{(\mu_{13})} V_{\mu_{13}}\bigr)
        + \Phi_{23}\bigl(B_{\ell(\mu_{23})+1}^{(\mu_{23})} V_{\mu_{23}}\bigr)
   \,\Bigr].
\end{equation}
\end{theorem}

\begin{proof}
We follow the May-15/May-20 pull-through, applied separately to each
of the four pieces of $\Rp_n V_\lambda$ from
Lemma~\ref{lem:Eplus-decomp}, then show two pieces vanish.

\medskip\textbf{Step 1 (pulling four $\Rp$-blocks through $\Phi$).}
Let $\Phi \in \{\Phi_{12}, \Phi_{13}, \Phi_{23}, \Phi_h\}$ and $\mu$ the
corresponding inner shape.

Decompose $\Rp_{n-1} = (T_{n-2}+1)\,U_1$ where $U_1 = (T_{n-3}+1)\cdots
(T_1+1) = \Rp_{n-2}^{(\mu)} \in \Hq(S_{n-2})$.  Since $T_1, \ldots,
T_{n-3}$ commute with $T_{n-1}$ (index gaps $\ge 2$), $U_1$ preserves
$E_{n-1}^+|_{V_\lambda}$.  By
Lemma~\ref{lem:Eplus-decomp}(i), $U_1 \Phi(v) = \Phi(\Rp_{n-2}^{(\mu)} v)$
for $v \in V_\mu$.  Hence
\[
   \Rp_{n-1}\Phi(V_\mu)
   \;=\; (T_{n-2}+1)\,\Phi\bigl(\Rp_{n-2}^{(\mu)} V_\mu\bigr).
\]

Iterating the same decomposition for $\Rp_{n-2}, \Rp_{n-3}, \Rp_{n-4}$,
using that each $U_k$ commutes with the previously-pulled $T_j$
factors (index gaps $\ge 2$):
\begin{equation}\label{eq:branchpull}
   \Rp_{n-4}\Rp_{n-3}\Rp_{n-2}\Rp_{n-1}\Phi(V_\mu)
   \;=\; (T_{n-5}+1)(T_{n-4}+1)(T_{n-3}+1)(T_{n-2}+1)\,
   \Phi\bigl(Q_\mu V_\mu\bigr),
\end{equation}
where
\[
   Q_\mu \;:=\; \Rp_{n-5}^{(\mu)}\, \Rp_{n-4}^{(\mu)}\, \Rp_{n-3}^{(\mu)}\, \Rp_{n-2}^{(\mu)}.
\]

\medskip\textbf{Step 2 (combining and applying Phase A).}
By Lemma~\ref{lem:Eplus-decomp}(ii) and Theorem~\ref{thm:phaseA},
\begin{align*}
   B \;=\; \Om V_\lambda
   &\;=\; \Rp_{n-4}\Rp_{n-3}\Rp_{n-2}\Rp_{n-1} \Rp_n V_\lambda \\
   &\;=\; \Rp_{n-4}\Rp_{n-3}\Rp_{n-2}\Rp_{n-1}
      \!\!\!\bigoplus_{X \in \{12,13,23,h\}}\!\!\!\Phi_X(V_{\mu_X}) \\
   &\;=\; (T_{n-5}{+}1)(T_{n-4}{+}1)(T_{n-3}{+}1)(T_{n-2}{+}1)\,
        \sum_{X}\Phi_X\bigl(I_X\bigr),
\end{align*}
where $I_X := Q_{\mu_X} V_{\mu_X}$.

\medskip\textbf{Step 3 (the $\mu_{12}$ piece vanishes).}
$\mu_{12} = (3, 2, 1^q)$ has $\ell(\mu_{12}) = q + 2 = n - 5$, hence
$\ell(\mu_{12}) + 1 = n - 4$.  The inner quadruple product is
\[
   I_{12} \;=\; \Rp_{n-5}^{(\mu_{12})}\,\bigl(\Rp_{n-4}^{(\mu_{12})}\,
   \Rp_{n-3}^{(\mu_{12})}\,\Rp_{n-2}^{(\mu_{12})} V_{\mu_{12}}\bigr)
   \;=\; \Rp_{n-5}^{(\mu_{12})}\,B^{(\mu_{12})}_{n-4} V_{\mu_{12}}.
\]
By Theorem~\ref{thm:32-prefix} applied to $\mu_{12} = (3, 2, 1^{q})$
with $q' = q \ge 3$ (joint hypothesis of the theorem and the present
proof; we will tighten to $q \ge 5$ below),
$B^{(\mu_{12})}_{n-4} V_{\mu_{12}} \subseteq E_1^-|_{V_{\mu_{12}}}$.
But $\Rp_{n-5}^{(\mu_{12})} = (T_{n-6}+1)(T_{n-7}+1)\cdots(T_1+1)$ has
$(T_1+1)$ as its \emph{rightmost} factor, which annihilates
$E_1^-|_{V_{\mu_{12}}}$.  Hence $I_{12} = 0$ and the $\mu_{12}$
contribution drops out.

\medskip\textbf{Step 4 (the hook piece vanishes).}
$\mu_h = (4, 1^{q+1})$ has $\ell(\mu_h) = q + 2 = n - 5$, hence
$\ell(\mu_h) + 1 = n - 4$.  Same argument as Step~3:
\[
   I_h \;=\; \Rp_{n-5}^{(\mu_h)}\,B^{(\mu_h)}_{n-4} V_{\mu_h}.
\]
By Theorem~\ref{thm:hook-E1minus} applied to $\mu_h = (4, 1^{q+1})$
with $k = 4$, $m = q + 1$ (hypothesis $m \ge \max(k, 2) = 4$, i.e.,
$q \ge 3$ $\checkmark$),
$B^{(\mu_h)}_{n-4} V_{\mu_h} \subseteq E_1^-|_{V_{\mu_h}}$.  The
rightmost $(T_1+1)$ of $\Rp_{n-5}^{(\mu_h)}$ kills this, so $I_h = 0$.

\medskip\textbf{Step 5 (the $\mu_{13}$ and $\mu_{23}$ pieces are
standard $B^{(\mu_X)}$).}  Both $\mu_{13} = (3, 3, 1^{q-1})$ and
$\mu_{23} = (4, 2, 1^{q-1})$ have $\ell(\mu_X) = q + 1 = n - 6$, hence
$\ell(\mu_X) + 1 = n - 5$.  The inner quadruple product runs from
$\Rp^{(\mu_X)}_{n-5}$ to $\Rp^{(\mu_X)}_{n-2}$, exactly the four blocks
of $\Om^{(\mu_X)}$.  Therefore
\[
   I_{13} \;=\; B^{(\mu_{13})}_{\ell(\mu_{13})+1} V_{\mu_{13}},
   \qquad
   I_{23} \;=\; B^{(\mu_{23})}_{\ell(\mu_{23})+1} V_{\mu_{23}}.
\]

Combining Steps 2--5 gives~\eqref{eq:reduction}.\qed
\end{proof}

\begin{remark}\label{rem:dim-add}
The dimensions match nicely:
\[
   \dim B \;=\; f^{(3, 2)} \;=\; 5,
   \qquad
   \dim B^{(\mu_{13})}_{\ell+1} \;=\; f^{(2,2)} \;=\; 2,
   \qquad
   \dim B^{(\mu_{23})}_{\ell+1} \;=\; f^{(3,1)} \;=\; 3,
\]
and $5 = 2 + 3$ confirms that both contributing Phi-pieces inject and
the four outer $(T_j+1)$ factors are injective on the image (computationally
verified at $q = 5$).
\end{remark}

\section{Main theorem}\label{sec:main}

\begin{theorem}[Main]\label{thm:main}
For $\lambda = (4, 3, 1^q) \vdash n = q + 7$ with $q \ge 5$,
\[
   B \;\subseteq\; \bigcap_{j=1}^{q-4} E_j^-\!\mid_{V_\lambda}.
\]
\end{theorem}

\begin{proof}
By Remark~\ref{rem:prefix-implies-Ej-}, it suffices to show that every
$T \in \supp(B)$ satisfies $p(T) \ge q - 3 = n - 10$.  We use the
two-branch reduction~\eqref{eq:reduction} together with prefix bounds
on the two inner $B^{(\mu_X)}_{\ell+1}$ pieces and
Lemma~\ref{lem:outer-preserve} on the four outer factors.

\medskip\textbf{Step 1 (prefix bound on the $\Phi_{13}$-image).}
$\mu_{13} = (3, 3, 1^{q-1})$ has size $n - 2$, $\ell(\mu_{13}) = q+1$,
and trailing-$1$-count $q' = q - 1$.  Since $q \ge 5$ we have
$q' \ge 4$, so Theorem~\ref{thm:33-prefix} applies, giving
\[
   p(T') \;\ge\; q' - 2 \;=\; q - 3
   \qquad \text{for every } T' \in \supp(B^{(\mu_{13})}_{\ell(\mu_{13})+1} V_{\mu_{13}}).
\]

Now $\Phi_{13}(v_{T'})$ is supported on the two SYTs of $\lambda$
obtained from $T'$ by placing $\{n-1, n\}$ at $\{c_1, c_3\} =
\{(1, 4), (q+2, 1)\}$ in the two orders.  Column-$1$ cells of
$\mu_{13}$ are at rows $1, 2, \ldots, q+1$ (i.e., $\ell(\mu_{13})$
rows).  The added cells $(1, 4)$ (column $4$, not $1$) and $(q+2, 1)$
(row $q+2$ at column $1$, with value $n - 1$ or $n \ne q + 2$) do not
extend the identity prefix.  Since $p(T') \le q + 1 = \ell(\mu_{13})$
and rows $1, \ldots, p(T')$ of column $1$ are unchanged, we get
$p(\Phi_{13}\text{-supp}) \ge p(T') \ge q - 3$.

\medskip\textbf{Step 2 (prefix bound on the $\Phi_{23}$-image).}
$\mu_{23} = (4, 2, 1^{q-1})$ has size $n - 2$, $\ell(\mu_{23}) = q+1$,
and trailing-$1$-count $q' = q - 1$.  Since $q \ge 5$ we have
$q' \ge 4$, so Theorem~\ref{thm:42-prefix} applies, giving
\[
   p(T') \;\ge\; q' - 2 \;=\; q - 3
   \qquad \text{for every } T' \in \supp(B^{(\mu_{23})}_{\ell(\mu_{23})+1} V_{\mu_{23}}).
\]
The Phi-embedding $\Phi_{23}$ adds $\{n-1, n\}$ at $\{c_2, c_3\} =
\{(2, 3), (q+2, 1)\}$: $(2, 3)$ is in column $3$ (not $1$), and
$(q+2, 1)$ is in row $q+2$ with value $n-1$ or $n \ne q+2$.  By the
same argument as Step~1, $p(\Phi_{23}\text{-supp}) \ge p(T') \ge q - 3$.

\medskip\textbf{Step 3 (outer factors preserve the prefix).}
The four outer factors in~\eqref{eq:reduction} are $(T_j{+}1)$ for
$j \in \{n-5, n-4, n-3, n-2\}$.  Each index satisfies $j \ge n - 5
= q + 2 > q - 2 = (q - 3) + 1$, so the hypothesis $j \ge P + 1$ of
Lemma~\ref{lem:outer-preserve} (with $P = q - 3$) is met.  Applying
the lemma four times in sequence, the prefix bound $\ge q - 3$ is
preserved.

\medskip\textbf{Conclusion.}  By Steps 1--3 and~\eqref{eq:reduction},
every $T \in \supp(B)$ has $p(T) \ge q - 3$.  By
Remark~\ref{rem:prefix-implies-Ej-}, $B \subseteq E_j^-$ for $1 \le
j \le q - 4 = \tau(\lambda)$.\qed
\end{proof}

\begin{remark}[Hypotheses]\label{rem:hyp}
The hypothesis $q \ge 5$ is the joint requirement of:
\begin{itemize}[topsep=0.2em,itemsep=0.1em]
\item Theorem~\ref{thm:33-prefix} at $q' = q - 1 \ge 4$, i.e., $q \ge 5$
(for $\mu_{13}$).
\item Theorem~\ref{thm:42-prefix} at $q' = q - 1 \ge 4$, i.e., $q \ge 5$
(for $\mu_{23}$).
\item Theorem~\ref{thm:32-prefix} at $q' = q \ge 3$, i.e., $q \ge 3$
(for $\mu_{12}$ vanishing in Step~3 of Theorem~\ref{thm:reduction}).
\item Theorem~\ref{thm:hook-E1minus} at $m = q + 1 \ge 4$, i.e., $q \ge 3$
(for $\mu_h$ vanishing in Step~4).
\end{itemize}
The threshold for the two recursive-input theorems ($q \ge 5$) is
tight: at $q = 4$ the prediction $\tau = q - 4 = 0$ is vacuous, and at
$q \le 3$ the prediction is negative.  Thus the hypothesis $q \ge 5$
exactly matches the first $q$ at which $\tau \ge 1$.
\end{remark}

\section{Computational verification}\label{sec:verify}

We verify Theorem~\ref{thm:main} and the structural decomposition
computationally at $q = 5$ (so $n = 12$, $\dim V_\lambda = 1232$).

\paragraph{Decomposition.}  At $q = 5$,
\begin{align*}
   f^{\mu_{12}} &= f^{(3,2,1^5)} = 160, &
   f^{\mu_{13}} &= f^{(3,3,1^4)} = 225, \\
   f^{\mu_{23}} &= f^{(4,2,1^4)} = 350, &
   f^{\mu_h} &= f^{(4,1^6)} = 84.
\end{align*}
Sum $= 819 = \dim \Rp_n V_\lambda = \dim E_{n-1}^+|_{V_\lambda}$,
consistent with Lemma~\ref{lem:Eplus-decomp}.

\paragraph{Vanishing.}  The two pulled-through inner pieces
$\Phi_{12}(I_{12})$ and $\Phi_h(I_h)$ contribute $0$ to $B$, leaving
the two pieces $\Phi_{13}(I_{13}) + \Phi_{23}(I_{23})$ with
$\dim I_{13} = 2$, $\dim I_{23} = 3$, summing to $5 = \dim B$.

\paragraph{$E_1^-$ containment at $q = 5$.}  Apply $\Om$ to five
independent random vectors in $V_\lambda$ (mod $P = 100\,003$ at $q = 11$)
and check $(T_1+1)\,\Om v = 0$ on each.  All five tests pass.  Probing
$\dim B$ via $\Om$ applied to the first $20$ basis vectors gives
rank $5$, matching $f^{(3,2)} = 5$.

\paragraph{Sharpness of the $q \ge 5$ hypothesis.}  At $q = 4$
(where $\tau = q - 4 = 0$ is vacuous), the same random-vector test
returns $164$ nonzero entries in $(T_1+1)\,\Om v$ on each of $5$
random vectors --- i.e., $B \not\subseteq E_1^-$ at $q = 4$.  The
boundary $q \ge 5$ for $B \subseteq E_1^-$ in the $\hatl = (4, 3)$
family is therefore sharp.

Scripts:
\texttt{\textasciitilde/projects/scratch/2026-05-14-extended-sign-kill-43/verify\_decomp.py,
verify\_E1minus\_q5.py, check\_q4\_sharp.py}.

\section{Updated table of proved hat\_lambda rows}\label{sec:table}

\begin{table}[h]
\centering
\renewcommand{\arraystretch}{1.15}
\caption{Status of the extended sign-kill conjecture by $\hatl$.  The
$\tau$ column lists the predicted threshold $\tau(\hatl) = q + 3 -
\hatl_1 - \hatl_2$.}
\begin{tabular}{l|c|l|l}
\toprule
$\hatl$ & predicted $\tau$ & status & paper \\
\midrule
$(2, 2)$ & $q - 1$ & proved $q \ge 2$           & May-14 fourth pass \\
$(3, 2)$ & $q - 2$ & proved $q \ge 3$           & May-14 fifth pass \\
$(4, 2)$ & $q - 3$ & proved $q \ge 4$ (gap closable) & May-14 sixth pass \\
$(3, 3)$ & $q - 3$ & proved $q \ge 4$           & May-14 seventh pass \\
$(4, 3)$ & $q - 4$ & \textbf{proved $q \ge 5$ (this paper)} & May-14 eighth pass \\
\bottomrule
\end{tabular}
\end{table}

\section{Discussion}\label{sec:discussion}

\subsection{The recursive ladder is consolidating}

This paper completes the second row of the recursive-ladder table:
every $\hatl$ with $\hatl_1 + \hatl_2 \le 7$ is now proved.  The next
natural target is $\hatl \in \{(5, 2), (4, 4), (5, 3)\}$, each with
threshold $\tau = q - 4$ or $\tau = q - 5$.

\subsection{Branch-vanishing as a structural feature}

Four pieces of $E_{n-1}^+$, two vanish, two contribute --- this is the
first case where \emph{two} branches drop out.  The pattern: a branch
$\mu_X$ vanishes precisely when $\ell(\mu_X) + 1 = n - 4$ (one less
than the standard ``$\ell + 1 = n - 5$'' alignment of $\mu_{13}, \mu_{23}$),
giving an extra inner $\Rp^{(\mu_X)}_{\ell(\mu_X)}$ factor whose
rightmost $(T_1+1)$ kills $B^{(\mu_X)}_{\ell+1} V_{\mu_X}$ via either
the inductive extended sign-kill (here, the $(3, 2)$ theorem at
$\mu_{12}$) or the hook column-$1$ theorem (at $\mu_h$).

This bookkeeping --- ``which $\mu_X$ have $\ell + 1 = n - 5$ vs.\
$n - 4$'' --- is shape-uniform once the corner-pair indexing is
fixed, and should generalize cleanly to deeper $\hatl$.

\subsection{Sharpness}

The sharpness direction ($B \not\subseteq E_{q-3}^-$ for $q \ge 5$)
would follow a structural reduction analogous to the seventh-pass
Sharpness Lemma (locating entry $n - 7$ in a witness SYT and arguing
via off-diagonal block + universal $T_\ell = q$).  We leave the
sharpness write-up for a follow-up note; the predicted threshold
matches the computational data at $q = 5$ (Table~1 of the conjecture
paper extended to $q = 5$: $B \not\subseteq E_2^-$ verified).

\subsection{Push status}

\textbf{69 unpushed commits} on \texttt{clio-vega/proofs} (after
today's commit).  PAT issue from May~6 persists --- please refresh.

\end{document}
